../cictikz/src/cictikz/data/tex/cictikz_preamble.tex

% Cartesian transmitter: build the symbol out of its real and imaginary
% parts rather than out of a phase and an amplitude.
%
% Two converters produce I and Q, two mixers multiply them by copies of
% the oscillator ninety degrees apart, and the sum is a signal with
% arbitrary phase and amplitude. No block here has to do anything hard
% except the amplifier, which has to be linear.
%
% Two deviations from the hand drawn original, both deliberate. The
% original drew the mixers with a plus sign; a mixer multiplies, so they
% are crosses here, and the place where the two branches meet gets the
% plus instead. And the original did not show the ninety degree shift on
% the oscillator, which is the one thing that makes this a Cartesian
% modulator rather than two copies of the same branch.
%
% The oscillator sits between the two branches rather than below them.
% Anywhere else and its wire to the far mixer has to cross a signal path,
% which in a block diagram reads as a connection that is not there. From
% the middle it fans out symmetrically and crosses nothing, and the
% ninety degree block lands in the branch it belongs to.

\begin{tikzpicture}[thick]
  %%%%%%%%%%%%%%%%%%%%%%%%%%%%%%%%%%%%%%%%%%%%%%%%%%%%%%%%%%%%%%%%%%%%%%
%% Shared pieces for the l04_afe signal flow graphs, l4_first_order and
%% l4_biquad.  The biquad is meant to read as the first-order graph with a
%% second integrator added, so the summing node and the 1/s block have to be
%% the same size and shape in both.  They are the same size in the original
%% artwork too.
%%
%% Names are prefixed and letters only, and computed coordinates are braced:
%% TikZ scans a coordinate for its closing parenthesis, so a '(' inside an
%% expression ends it early.
%%%%%%%%%%%%%%%%%%%%%%%%%%%%%%%%%%%%%%%%%%%%%%%%%%%%%%%%%%%%%%%%%%%%%%

\newcommand{\sfgR}{0.25}    % summing node radius
\newcommand{\sfgBw}{1.0}    % half width of the 1/s block
\newcommand{\sfgBh}{0.65}   % half height of the 1/s block

% A summing node at (#1,#2).
\newcommand{\sfgSum}[2]{%
  \draw (#1,#2) circle (\sfgR);
  \draw ({#1-0.17},#2) -- ({#1+0.17},#2);
  \draw (#1,{#2-0.17}) -- (#1,{#2+0.17});
}

% A block centred at (#1,#2) carrying #3.
\newcommand{\sfgBox}[3]{%
  \draw ({#1-\sfgBw},{#2-\sfgBh}) rectangle ({#1+\sfgBw},{#2+\sfgBh});
  \node at (#1,#2) {#3};
}

\newcommand{\sfgDot}[2]{\fill (#1,#2) circle (0.07);}


  \def\yi{1.3}
  \def\yq{-1.3}
  \def\xm{-1.0}          % the mixer column, and the oscillator trunk
  \def\xs{1.7}           % the summing node

  % the two converters
  \draw (-4.6,{\yi-0.5}) rectangle (-3.0,{\yi+0.5});
  \node at (-3.8,\yi) {DAC};
  \draw (-4.6,{\yq-0.5}) rectangle (-3.0,{\yq+0.5});
  \node at (-3.8,\yq) {DAC};
  \node[anchor=east] at (-5.5,\yi) {$I$};
  \node[anchor=east] at (-5.5,\yq) {$Q$};
  \draw[->] (-5.4,\yi) -- (-4.6,\yi);
  \draw[->] (-5.4,\yq) -- (-4.6,\yq);

  % mixers, fed from the converters
  \sfgMix{\xm}{\yi}
  \sfgMix{\xm}{\yq}
  \draw (-3.0,\yi) -- ({\xm-\sfgR},\yi);
  \draw (-3.0,\yq) -- ({\xm-\sfgR},\yq);

  % the oscillator, in the gap between the branches, feeding the I mixer
  % from below and the Q mixer from above
  %- same footprint as the two converters, so the three sources read as
  %- one column
  \draw (-4.6,-0.5) rectangle (-3.0,0.5);
  \node at (-3.8,0) {LO};
  \draw (-3.0,0) -- (\xm,0);
  \sfgDot{\xm}{0}
  \draw (\xm,0) -- (\xm,{\yi-\sfgR});
  %- the shift belongs on Q: s(t) = I cos(wt) - Q sin(wt), and sine is
  %- cosine ninety degrees late
  \draw ({\xm-0.45},-0.42) rectangle ({\xm+0.45},-0.98);
  \node[scale=0.8] at (\xm,-0.7) {$90^\circ$};
  \draw (\xm,0) -- (\xm,-0.42);
  \draw (\xm,-0.98) -- (\xm,{\yq+\sfgR});

  % sum the two branches
  \sfgSum{\xs}{0}
  \draw ({\xm+\sfgR},\yi) -- (\xs,\yi) -- (\xs,\sfgR);
  \draw ({\xm+\sfgR},\yq) -- (\xs,\yq) -- (\xs,{-\sfgR});

  \draw ({\xs+\sfgR},0) -- (2.9,0);
  \sfgAmp{3.8}{0}{PA}
  \sfgAntenna{4.8}{0}
  \draw (3.8,0) -- (4.8,0);
\end{tikzpicture}
\end{document}
